\section{Introducción}

El electrocardiograma (ECG) es una de las mediciones más comunes en medicina: con unos pocos electrodos colocados sobre la piel se puede registrar la actividad eléctrica del corazón de forma simple y no invasiva.
Sin embargo, medir esta señal no es tan directo como parece.
Su amplitud es muy chica, del orden de los milivolts, y suele llegar mezclada con interferencias bastante más grandes, como el ruido de la red eléctrica de 50~Hz, además de variaciones lentas de la línea de base causadas por la respiración, el movimiento o el contacto entre el electrodo y la piel.
Por eso, antes de poder visualizar un ECG es necesario amplificar la señal y filtrarla de manera adecuada.

La primera etapa de este proceso suele ser un amplificador de instrumentación, ya que tiene una impedancia de entrada muy alta, una ganancia fácil de ajustar y un buen rechazo de las señales de modo común.
En este trabajo práctico se armó en protoboard un canal básico de ECG usando el amplificador de instrumentación AD620 junto con los dos amplificadores operacionales del integrado TL072.
El circuito está formado por una etapa de preamplificación, una etapa de filtrado y una etapa de ganancia adicional.

El objetivo principal fue entender cómo funciona un amplificador de instrumentación en una aplicación real y comprobar en la práctica el funcionamiento del canal de ECG.
Para eso, se probaron distintos valores de la resistencia de ganancia $R_G$ del AD620 para ver cómo cambiaba la señal obtenida, y se analizó la etapa de filtrado para entender cómo las frecuencias de corte afectan la forma del ECG.
A partir de esto, se propusieron posibles mejoras en el filtrado, la ganancia y la alimentación del circuito.